\wordchapter{7}{Conclusion}

This thesis examines the design of the government support to the crop disaster insurance and the production patterns when the support was introduced in South Korea. The institutional assessment answers three questions: Premium support and access and distribution of premium support; Administrative expenditure and its output; Reinsurance and risk allocation and public financial exposure.

Across six annual crops, the mean production contrast is −0.002 log points, with an exploratory 95 percent interval of [−0.318, 0.315]. Estimates vary with reference periods, comparison crops, and weights, and area and production differ in pre-pilot growth. For municipal rice area, the estimate is −0.021 log points, with a 95 percent interval of [−0.064, 0.022], and changes with sample composition. The rice cohorts have closer observed pre-pilot paths, but the analysis covers only two years after introduction. Together, these results leave substantial uncertainty about production responses.

The results support evaluating production alongside the other purposes of public support. For premium subsidies, this requires evidence on who receives assistance and who obtains protection. For administrative support, it requires relating expenditure to comparable measures of service and cost. For reinsurance, it requires comparing risk retention and public financing under alternative loss-sharing terms. Production responses form one part of this assessment; they do not alone determine the appropriate level or design of government support.
